\documentclass[10pt,a4paper]{article}

% ---------- Geometry ----------
\usepackage[a4paper,left=2.0cm,right=2.0cm,top=1.9cm,bottom=1.9cm]{geometry}
\setlength{\parindent}{0pt}
\setlength{\parskip}{0.35em}

% ---------- Encoding / fonts ----------
\usepackage[utf8]{inputenc}
\usepackage[T1]{fontenc}
\usepackage[english]{babel}
\usepackage{microtype}
\usepackage{mathptmx} % Times-like text+math, compact and print-friendly

% ---------- Math / physics ----------
\usepackage{amsmath}
\usepackage{amssymb}
\usepackage{mathtools}
\usepackage{braket}
\usepackage[separate-uncertainty=true]{siunitx}
\usepackage{isotope}

% ---------- Figures ----------
\usepackage{graphicx}
\usepackage{caption}
\usepackage{svg}
\captionsetup{font=small,labelfont=bf,skip=4pt}

% ---------- Misc ----------
\usepackage[hidelinks]{hyperref}
\usepackage{enumitem}
\setlist{nosep,leftmargin=1.4em}
\usepackage{titlesec}
\titlespacing*{\section}{0pt}{0.7em}{0.35em}
\titleformat{\section}{\normalfont\bfseries\large}{\thesection}{0.6em}{}
\pagestyle{empty}

\begin{document}

% ============ Title block ============
\begin{center}
	{\LARGE\bfseries An Atom--Photon Quantum Gate Using an Atomic Clock Qubit and a Birefringent Cavity}\\[0.6em]
	{\large Leart Zuka}\\[0.15em]
	{\small M.Sc. Condensed Matter Physics, Technical University of Munich\\
	Thesis carried out at the Max Planck Institute of Quantum Optics, Garching\\
	Supervisor: Prof.\ Dr.\ Gerhard Rempe}
\end{center}

\vspace{0.3em}

% ============ Motivation ============
\section*{Motivation}

Quantum networks aim to link spatially separated quantum processors into a single, more powerful and more secure architecture by distributing entanglement between stationary qubits via flying qubits (in this case photons). Each network node must be able to perform three elementary tasks: store the incoming information of a photonic qubic onto a stationary qubits (storage), generate entanglement between the stationary and photonic qubit, and \emph{process} quantum information locally through a deterministic quantum logic gate. The \textit{CavityX} platform at the Max Planck Institute of Quantum Optics aims at realizing such a node using a single \isotope[87]{Rb} atom coupled to a birefringent fiber Fabry--P\'erot cavity. It has previously demonstrated photon-to-atom state storage (\SI{94.7}{\percent} fidelity) and the generation of heralded atom--photon entanglement (\SI{87}{\percent} fidelity) -- the first two node primitives. This thesis supplies the missing third primitive by realizing a deterministic atom--photon quantum gate, following the protocol proposed by Duan and Kimble [1].

A further design goal is robustness: qubits encoded in magnetically sensitive atomic states dephase due to ambient magnetic field fluctuations. This work instead encodes the atomic qubit in the magnetic-field-insensitive \emph{clock states} of the \isotope[87]{Rb} $5S_{1/2}$ ground manifold, which is possible because -- unlike in prior realizations of the Duan--Kimble protocol [2] -- the required polarization selectivity is provided by the cavity's birefringence rather than by Zeeman splitting.

% ============ Working principle ============
\section*{Working Principle}

The cavity's two linear polarization eigenmodes are split by $\Delta_{\text{biref}} \approx 2\pi\times\SI{505}{\mega\hertz}$. Only the $\pi$-polarized mode is resonant with the atomic $\pi$-transition. The orthogonal $V$-mode on the other hand is far detuned and reflects off the empty cavity regardless of the atomic state. A $\pi$-polarized photon reflecting off the atom-cavity system therefore picks up an atom-state-dependent phase: a $\pi$ phase shift if the atom is in the uncoupled clock state $\ket{0}=\ket{F{=}1,m_F{=}0}$, and no phase shift if the atom is in the coupled clock state $\ket{1}=\ket{F{=}2,m_F{=}0}$ (Fig.~1). Because the atom is the control qubit and the photon polarization is the target qubit, this state-dependent reflection directly implements a controlled-phase-flip (CPF) gate in the photonic $\{\ket{\pi},\ket{V}\}$ basis. Preparing the photon instead in a superposition of $\pi$ and $V$ -- realized experimentally using circular polarization, $\ket{R/L} = (\ket{V}\pm\ket{\pi})/\sqrt{2}$ -- converts the same interaction into a controlled-NOT (CNOT) gate.

\begin{figure}[!htb]
	\centering
	\includesvg[inkscapelatex=false,width=\linewidth]{Figures/5_atom_photon_gate/nms_reflection_phase_horizontal.svg}
	\caption{Normal-mode spectroscopy of the atom-cavity system. (a) Measured reflectivity $|r|^2$ versus probe detuning $\Delta$ for the atom in the uncoupled state $\ket{0}$ (red) and coupled state $\ket{1}$ (blue), with fits to the analytical model above; the coupled spectrum shows the vacuum Rabi splitting from which $g$ is extracted. (b) Phase $\arg(r)$ computed from the same fitted model, showing the predicted $\pi$-phase-shift asymmetry between the two atomic states.}
\end{figure}

Because the $V$-mode reflects with near-unity amplitude while the $\pi$-mode reflectivity is reduced by intracavity losses, an array of Brewster windows was installed in the input path to attenuate the $V$-polarized component and balance the two reflection amplitudes, maximizing the CNOT gate overlap with the ideal truth table.

% ============ Results ============
\section*{Experimental Results}

The atom--cavity coupling strength was determined from normal-mode spectroscopy to be $g = 2\pi\times(25.9\pm0.5)\,\si{\mega\hertz}$, consistent with a master-equation simulation used to identify the optimal operating point: a mean input photon number of $\bar{n}\approx0.1$ and seven Brewster windows ($T_V\approx\SI{24}{\percent}$). Operating at this point, the gate truth tables were measured by preparing the atom and photon in each of the four computational basis-state combinations and reading out the joint output state (Fig.~2). The measured truth tables overlap with the ideal CPF and CNOT gates by
\[
	\mathcal{O}_{\text{CPF}} = (92.41\pm0.52)\,\%, \qquad \mathcal{O}_{\text{CNOT}} = (88.15\pm0.86)\,\%,
\]
in good agreement with simulation. The dominant infidelity contributions are the residual asymmetry between the coupled and uncoupled $\pi$-mode reflectivities ($\sim\SI{5}{\percent}$, set by the current cavity cooperativity) and imperfect preparation of the coupled atomic clock state ($\sim\SI{6}{\percent}$).

\begin{figure}[!htb]
	\centering
	\includesvg[width=\linewidth]{Figures/5_atom_photon_gate/gate_truth_tables_2x2.svg}
	\caption{Measured gate truth tables at $\bar{n}\approx0.1$: (c) CPF gate in the linear $\{\ket{V},\ket{\pi}\}$ photon basis; (d) CNOT gate in the circular $\{\ket{R},\ket{L}\}$ photon basis. Diagonal elements give the probability of the correct output state for each input combination.}
\end{figure}

To confirm the coherent, entangling action of the gate rather than merely its classical input--output behavior, an atom--photon Bell state $\ket{\Phi^+}$ was generated from a separable input and characterized via quantum state tomography, yielding a fidelity of $(65.90\pm1.73)\,\%$ -- exceeding the \SI{50}{\percent} threshold required to certify entanglement, and thus confirming genuine atom--photon entanglement generated by the gate.

\begin{figure}[!htb]
	\centering
	\includesvg[width=\linewidth]{Figures/5_atom_photon_gate/bell_state_L.svg}
	\caption{Magnitude of the reconstructed atom--photon density matrix $|\rho_{ij}|$ for the generated Bell state $\ket{\Phi^+}$, obtained via quantum state tomography. Population and coherence are concentrated on the $\{\ket{0,R},\ket{1,L}\}$ subspace, as expected for $\ket{\Phi^+}$, with residual leakage into the wrong-parity states.}
\end{figure}

% ============ Outlook ============
\section*{Conclusion and Outlook}

This thesis experimentally realizes a deterministic atom--photon CPF/CNOT gate on the CavityX platform, completing the set of primitives -- storage, entanglement generation, and now processing -- required for a functional quantum-network node, while encoding the atomic qubit in magnetic-field-insensitive clock states for improved robustness. The next steps would be to systematically optimize the atomic state-preparation fidelity to improve the Bell-state fidelity, followed by inter-node protocols such as quantum state teleportation over the existing \SI{24}{\kilo\metre} fiber link between MPQ Garching and LMU Munich [3]. Longer term, the same gate mechanism enables a photon--photon gate mediated by the single atom, and -- by combining it with the platform's quantum-memory capability -- the heralded generation of entanglement between temporally separated photons.

\vspace{0.3em}
{\small
	\textbf{References}
	\begin{enumerate}[label={[\arabic*]},leftmargin=2.0em,itemsep=0pt,topsep=2pt]
		\item L.-M. Duan, H. J. Kimble, \emph{Phys. Rev. Lett.} \textbf{92}, 127902 (2004).
		\item A. Reiserer, N. Kalb, G. Rempe, S. Ritter, \emph{Nature} \textbf{508}, 237 (2014).
		\item M. Büki \emph{et al.}, \emph{Phys. Rev. Lett.} \textbf{137}, 090803 (2026)
	\end{enumerate}
}

\end{document}
